\begin{thebibliography}{99}

\bibitem{xie1}
Ying Xie.
\emph{What the Machine Invented: Mining New Mathematical Instruments from OpenAI's October 2026 Release of AI-Generated Mathematics}.
Preprint, version 1, October 12, 2026.
\href{https://doi.org/10.5281/zenodo.23319145}{doi:10.5281/zenodo.23319145}.

\bibitem{xie2}
Ying Xie.
\emph{From Mathematical Instruments to Reusable Reasoning: A Second Phase of Extraction from AI-Generated Mathematics}.
Preprint, version 1, October 21, 2026.
\href{https://doi.org/10.5281/zenodo.23523260}{doi:10.5281/zenodo.23523260}.

\bibitem{leiden}
Jarod Alper, Michael Barany, Alain Chavarri Villarello, Sander Dahmen, Walter Dean, Karthik Ganapathy, Michael Harris, David Holmes, Mateja Jamnik, Steven Kelk, Bryna Kra, Ursula Martin, Bartosz Naskr\k{e}cki, Rodrigo Ochigame, Jim Portegies, and Johannes Schmitt.
\emph{Leiden Declaration on Artificial Intelligence and Mathematics}.
Zenodo, June 2, 2026.
\href{https://doi.org/10.5281/zenodo.20302944}{doi:10.5281/zenodo.20302944}.

\bibitem{polya}
George P\'olya.
\emph{How to Solve It: A New Aspect of Mathematical Method}.
Expanded Princeton Science Library edition, with a foreword by John H. Conway. Princeton University Press, 2004. Originally published in 1945.
\href{https://doi.org/10.1515/9781400828678}{doi:10.1515/9781400828678}.

\bibitem{lakatos}
Imre Lakatos.
\emph{Proofs and Refutations: The Logic of Mathematical Discovery}.
Edited by John Worrall and Elie Zahar. Cambridge University Press, 1976.
\href{https://doi.org/10.1017/CBO9781139171472}{doi:10.1017/CBO9781139171472}.

\bibitem{bundy}
Alan Bundy.
The use of explicit plans to guide inductive proofs.
In Ewing Lusk and Ross Overbeek, editors, \emph{9th International Conference on Automated Deduction: Argonne, Illinois, USA, May 23--26, 1988. Proceedings},
Lecture Notes in Computer Science, vol.\ 310, pp.\ 111--120. Springer, 1988.
\href{https://doi.org/10.1007/BFb0012826}{doi:10.1007/BFb0012826}.

\bibitem{htps}
Guillaume Lample, Marie-Anne Lachaux, Thibaut Lavril, Xavier Martinet, Amaury Hayat, Gabriel Ebner, Aur\'elien Rodriguez, and Timoth\'ee Lacroix.
HyperTree Proof Search for Neural Theorem Proving.
In \emph{Advances in Neural Information Processing Systems 35 (NeurIPS 2022)}, pp.\ 26337--26349. Curran Associates, Inc., 2022.
\href{https://doi.org/10.52202/068431-1910}{doi:10.52202/068431-1910}.

\bibitem{mkm}
Michael Kohlhase, Florian Rabe, and Vyacheslav Zholudev.
Towards MKM in the Large: Modular Representation and Scalable Software Architecture.
In Serge Autexier, Jacques Calmet, David Delahaye, Patrick D. F. Ion, Laurence Rideau, Renaud Rioboo, and Alan P. Sexton, editors, \emph{Intelligent Computer Mathematics: 10th International Conference, AISC 2010, 17th Symposium, Calculemus 2010, and 9th International Conference, MKM 2010, Paris, France, July 5--10, 2010. Proceedings}, Lecture Notes in Computer Science, vol.\ 6167, pp.\ 370--384. Springer, 2010.
\href{https://doi.org/10.1007/978-3-642-14128-7_32}{doi:10.1007/978-3-642-14128-7\_32}.

\bibitem{mathlib}
The mathlib Community.
The Lean Mathematical Library.
In Jasmin Blanchette and C\u{a}t\u{a}lin Hri\c{t}cu, editors, \emph{Proceedings of the 9th ACM SIGPLAN International Conference on Certified Programs and Proofs (CPP 2020)}, New Orleans, LA, USA, January 20--21, 2020, pp.\ 367--381. ACM, 2020.
\href{https://doi.org/10.1145/3372885.3373824}{doi:10.1145/3372885.3373824}.

\bibitem{polymath}
Timothy Gowers and Michael Nielsen.
Massively collaborative mathematics.
\emph{Nature} 461(7266), 879--881, 2009.
\href{https://doi.org/10.1038/461879a}{doi:10.1038/461879a}.

\bibitem{prov}
Paul Groth and Luc Moreau, editors.
\emph{PROV-Overview: An Overview of the PROV Family of Documents}.
W3C Working Group Note, April 30, 2013.
\url{https://www.w3.org/TR/2013/NOTE-prov-overview-20130430/}.

\bibitem{funsearch}
Bernardino Romera-Paredes, Mohammadamin Barekatain, Alexander Novikov, Matej Balog, M. Pawan Kumar, Emilien Dupont, Francisco J. R. Ruiz, Jordan S. Ellenberg, Pengming Wang, Omar Fawzi, Pushmeet Kohli, and Alhussein Fawzi.
Mathematical discoveries from program search with large language models.
\emph{Nature} 625(7995), 468--475, 2024.
\href{https://doi.org/10.1038/s41586-023-06924-6}{doi:10.1038/s41586-023-06924-6}.

\bibitem{aiscientist}
Chris Lu, Cong Lu, Robert Tjarko Lange, Jakob Foerster, Jeff Clune, and David Ha.
\emph{The AI Scientist: Towards Fully Automated Open-Ended Scientific Discovery}.
\href{https://arxiv.org/abs/2408.06292v3}{arXiv:2408.06292v3} [cs.AI], September 1, 2024.
\href{https://doi.org/10.48550/arXiv.2408.06292}{doi:10.48550/arXiv.2408.06292}.

\bibitem{valiant}
Leslie G. Valiant.
Completeness classes in algebra.
In \emph{Proceedings of the Eleventh Annual ACM Symposium on Theory of Computing}, Atlanta, GA, USA, April 30--May 2, 1979, pp.\ 249--261. ACM, 1979.
\href{https://doi.org/10.1145/800135.804419}{doi:10.1145/800135.804419}.

\bibitem{meshkat}
Nicolette Meshkat and Seth Sullivant.
Identifiable reparametrizations of linear compartment models.
\emph{Journal of Symbolic Computation} 63, 46--67, 2014.
\href{https://doi.org/10.1016/j.jsc.2013.11.002}{doi:10.1016/j.jsc.2013.11.002}.


\bibitem{landsbergmanivel}
J. M. Landsberg and Laurent Manivel.
On the ideals of secant varieties of Segre varieties.
\emph{Foundations of Computational Mathematics} 4(4), 397--422, 2004.
\href{https://doi.org/10.1007/s10208-003-0115-9}{doi:10.1007/s10208-003-0115-9}.
\href{https://arxiv.org/abs/math/0311388v1}{arXiv:math/0311388v1}.

\bibitem{raicu}
Claudiu Raicu.
Secant varieties of Segre--Veronese varieties.
\emph{Algebra \& Number Theory} 6(8), 1817--1868, 2012.
\href{https://doi.org/10.2140/ant.2012.6.1817}{doi:10.2140/ant.2012.6.1817}.
\href{https://arxiv.org/abs/1011.5867v2}{arXiv:1011.5867v2}.

\bibitem{ballicobernardi}
Edoardo Ballico and Alessandra Bernardi.
Tensor ranks on tangent developable of Segre varieties.
\emph{Linear and Multilinear Algebra} 61(7), 881--894, 2013.
\href{https://doi.org/10.1080/03081087.2012.716430}{doi:10.1080/03081087.2012.716430}.
\href{https://arxiv.org/abs/1210.7976v2}{arXiv:1210.7976v2}.

\bibitem{aliabadifriedland}
Mohsen Aliabadi and Shmuel Friedland.
On the complexity of finding tensor ranks.
\emph{Communications on Applied Mathematics and Computation} 3(2), 281--289, 2021.
\href{https://doi.org/10.1007/s42967-020-00103-4}{doi:10.1007/s42967-020-00103-4}.
\href{https://arxiv.org/abs/2002.07151v2}{arXiv:2002.07151v2}.

\bibitem{galesi}
Nicola Galesi, Joshua A. Grochow, Toniann Pitassi, and Adrian She.
On the Algebraic Proof Complexity of Tensor Isomorphism.
In \emph{38th Computational Complexity Conference (CCC 2023)},
Leibniz International Proceedings in Informatics, vol.\ 264, pp.\ 4:1--4:40. Schloss Dagstuhl--Leibniz-Zentrum f\"ur Informatik, 2023.
\href{https://doi.org/10.4230/LIPIcs.CCC.2023.4}{doi:10.4230/LIPIcs.CCC.2023.4}.

\bibitem{buczynskilandsberg}
Jaros{\l}aw Buczy\'nski and J. M. Landsberg.
Ranks of tensors and a generalization of secant varieties.
\emph{Linear Algebra and its Applications} 438(2), 668--689, 2013.
\href{https://doi.org/10.1016/j.laa.2012.05.001}{doi:10.1016/j.laa.2012.05.001}.

\end{thebibliography}
